\begin{textsample}{POS Dim 5 – vem\_ed\_24 – Score 12.00 – vem\_ed\_24/t128.txt}  \label{ex:f5_pos_012}
IVES Cesu 2024: Estratégias de colaboração internacional Sandra Valencia, da Universidad Católica de Manizales (Colômbia), \textbf{comentou} as estratégias de colaboração internacional na educação superior por meio de projetos COIL.
Apontou os mecanismos, \textbf{modelos} e formas para fazer Intercâmbios Virtuais que já estão institucionalizados internacionalmente.
Ressalvou que \textbf{poucas} oportunidades de formação profissional em cooperação internacional \textbf{dificultam} a compreensão dessas práticas.
Entre os principais desafios, estão a resistência, que alguns professores têm, em relação à aprendizagem de idiomas estrangeiros e o acesso \textbf{limitado} à internet, especialmente em \textbf{universidades} no interior do país – questão que não é exclusiva da Colômbia, mas \textbf{presente} em diversas paragens do Sul Global.
Sandra Valencia também frisou a importância de relações mais equitativas nos Intercâmbios Virtuais, entre países do Norte e do Sul Global.
Explanou tópicos como \textbf{compromisso} global, empoderamento pessoal e institucional, experiências globais e interações colaborativas on-line em equipes internacionais.
Mencionou o projeto “COIL for staff”, envolvendo funcionários administrativos de sua instituição e da Universidad CEU Cardenal Herrera (Espanha).
Flexibilidade, \textbf{adaptabilidade} e comunicação: eis as qualidades principais de um professor envolvido em projetos COIL. “Flexibilidade não só na gestão de tempo, mas nas ideias.
Pensar na colaboração, abrir para outros contextos e possibilidades e proporcionar ideias para seu parceiro a partir de suas forças e conhecendo as próprias fraquezas”.
Conexões latino-americanas Brenda García preside a Red LatAM COIL, fundada na pandemia para fortalecer projetos COIL entre \textbf{universidades} da América Latina e instituições de todo o mundo.
São três os objetivos principais da \textbf{rede}: 1.
Promover COIL entre a América Latina e o resto do mundo; 2.
Promover prática e investigação sobre o tema, em uma conferência anual on-line, que trata sobre Internacionalização em Casa, Internacionalização do Ensino Superior e do Currículo (\textbf{ver} agenda de eventos do segundo semestre); 3.
Expandir os benefícios da metodologia COIL como estratégia de internacionalização do \textbf{currículo} na educação superior por meio da colaboração com outras \textbf{redes} e organizações. “COIL pode ser uma das estratégias mais importantes para a internacionalização, oferecendo oportunidades para colaborar, cocriar e pesquisar”, acredita Brenda. “Meu maior sonho, desde antes de ocupar este honroso lugar que é a presidência da Red LatAM COIL, é empoderar estudantes e instituições da América Latina”.
Segundo a presidente da Red LatAM COIL, “ainda falta muita pesquisa para nos posicionar internacionalmente, mas, por meio dos projetos COIL, a internacionalização abriu os olhos do mundo para a América Latina”.
Ela acredita nos projetos COIL como uma alternativa permanente e sustentável para desenvolver competências interculturais entre os estudantes.

% matched lemmas: adaptabilidade, comentar, compromisso, currículo, dificultar, limitar, modelo, pouco, presente, rede, universidade, ver
\end{textsample}
